\documentclass[a4paper,12pt]{article}
\usepackage[T2A]{fontenc}
\usepackage[utf8]{inputenc}
\usepackage[russian]{babel}
\usepackage{amsmath,amssymb,mathtools}
\usepackage{helvet}
\renewcommand{\familydefault}{\sfdefault}
\usepackage{icomma}
\usepackage[margin=22mm,headheight=15pt,footskip=13mm]{geometry}
\usepackage{microtype}
\usepackage{booktabs,tabularx,array}
\usepackage{enumitem}
\usepackage{parskip}
\usepackage{fancyhdr}
\usepackage{tikz}
\usepackage{xcolor}
\usepackage[colorlinks=true,linkcolor=teal!55!black]{hyperref}
\definecolor{accent}{HTML}{087B83}
\definecolor{ink}{HTML}{252B30}
\definecolor{muted}{HTML}{67727B}
\definecolor{errorcolor}{HTML}{B23434}
\color{ink}
\setlength{\parindent}{0pt}
\setlength{\parskip}{7pt}
\setlist[enumerate]{itemsep=4pt,topsep=4pt,leftmargin=*}
\pagestyle{fancy}
\fancyhf{}
\fancyhead[L]{\small\textcolor{muted}{Математика · ЕГЭ, профильный уровень}}
\fancyhead[R]{\small\textcolor{muted}{4 ноября 2026 г.}}
\fancyfoot[C]{\scriptsize Лёвин Артём Александрович · эксклюзивно для Николь Саркисьянц}
\fancyfoot[R]{\scriptsize\thepage}
\renewcommand{\headrulewidth}{0pt}
\renewcommand{\footrulewidth}{0pt}
\newcommand{\keyidea}{\textcolor{accent}{\textbf{Ключевая идея.}}\ }
\newcommand{\wrong}{\textcolor{errorcolor}{\textbf{Ошибка.}}\ }
\newcommand{\result}{\textcolor{accent}{\textbf{Итог.}}\ }
\newcommand{\taskstart}[2]{\clearpage\phantomsection\label{task:#1}\section*{Задание № #1}\textcolor{accent}{#2}\par\vspace{2mm}}
\newcommand{\field}[1]{\par\smallskip\textbf{#1}\par\vspace{-3mm}}
\begin{document}
\thispagestyle{empty}
\begin{center}
\vspace*{31mm}
{\large\textcolor{accent}{МАТЕМАТИКА}}\par\vspace{22mm}
{\LARGE Ревью домашней работы}\par\vspace{7mm}
{\Large Профильный ЕГЭ: вариант с развёрнутыми решениями}\par\vspace{25mm}
{\large Николь Саркисьянц}\par\vspace{17mm}
{4 ноября 2026 г.}
\vfill
{\small Лёвин Артём Александрович}\par
{\small эксклюзивно для Николь Саркисьянц}\par\vspace{13mm}
\begin{tikzpicture}
\draw[accent,line width=.9pt] (0,0) .. controls (1.15,.55) and (2.2,-.55) .. (3.4,0) .. controls (4.5,.50) and (5.5,-.45) .. (6.6,0) .. controls (7.6,.4) and (8.7,-.45) .. (9.6,0);
\end{tikzpicture}
\end{center}
\clearpage
\section*{Сводная таблица}
В таблице представлены только задания, требующие исправления или завершения. Номер задания ведёт к подробному разбору.
\par\medskip
{\small\renewcommand{\arraystretch}{1.4}
\begin{tabularx}{\linewidth}{@{}>{\centering\arraybackslash}p{10mm} >{\raggedright\arraybackslash}p{28mm} X >{\raggedright\arraybackslash}p{24mm} >{\raggedright\arraybackslash}p{32mm}@{}}
\toprule
№ & Тема & Суть затруднения & Ваш ответ & Верный ответ \\
\midrule
\hyperref[task:3]{3} & Объёмы тел & Решение отсутствует & Нет & $76{,}5$ \\
\hyperref[task:5]{5} & Вероятность & Решение отсутствует & Нет & $0{,}15$ \\
\hyperref[task:6]{6} & Дискретная величина & Решение отсутствует & Нет & $0{,}55$ \\
\hyperref[task:9]{9} & Производная & Неверно подсчитаны точки максимума & $5$ & $2$ \\
\hyperref[task:11]{11} & Задача на движение & Время записано как скорость, делённая на путь & $8$ км & $390$ км \\
\hyperref[task:13]{13} & Кредит & Решение отсутствует & Нет & $1\,594\,200$ руб. \\
\hyperref[task:14]{14} & Тригонометрия & Ошибка в формуле приведения и работе с корнями & Не завершено & $-4\pi$, $-3\pi$, $-\frac{8\pi}{3}$ \\
\bottomrule
\end{tabularx}}
\par\medskip
\textcolor{muted}{Примечание.} При отсутствии письменного решения невозможно установить ошибочный промежуточный шаг. Для таких заданий приведён полный эталонный маршрут.
\clearpage
\thispagestyle{empty}
\vspace*{.32\textheight}
\begin{center}
{\LARGE\textcolor{accent}{Разбор ошибок и пропусков}}\par\vspace{9mm}
{\large Последний верный шаг · причина · исправление · проверка}
\end{center}

\taskstart{3}{Стереометрия · объём цилиндра}
\field{Условие}
Объём первого цилиндра равен $102$. У второго цилиндра высота в $3$ раза больше, а радиус основания в $2$ раза меньше. Найдите объём второго цилиндра.
\field{Ваше решение и ответ}
На листе отсутствуют.
\field{Где возникает затруднение}
Письменных преобразований нет, поэтому указать последний правильный и первый неправильный шаг невозможно. Для решения нужно учесть квадрат радиуса в формуле объёма.
\field{Ключевая идея}
\keyidea Объём цилиндра равен $V=\pi r^2h$. Изменение радиуса в $k$ раз изменяет площадь основания в $k^2$ раз.
\field{Полное правильное решение}
Пусть у первого цилиндра радиус $r$, высота $h$, объём $V_1=102$. Тогда
\[V_1=\pi r^2h.\]
У второго цилиндра радиус $r/2$, высота $3h$. Следовательно,
\begin{align*}
V_2&=\pi\left(\frac r2\right)^2\cdot 3h\\
&=\pi\frac{r^2}{4}\cdot3h
=\frac34\pi r^2h
=\frac34 V_1\\
&=\frac{3\cdot102}{4}=76{,}5.
\end{align*}
\field{Проверка}
Высота увеличилась в $3$ раза, площадь основания уменьшилась в $4$ раза; объём должен составить $3/4$ исходного. Полученное число этому соответствует.
\result $V_2=76{,}5$.
\field{Задача на закрепление}
Объём первого цилиндра равен $80$ см$^3$. У второго высота в $5$ раз больше, радиус основания в $2$ раза меньше. Найдите объём второго цилиндра.

\taskstart{5}{Теория вероятностей · пересечение событий}
\field{Условие}
В двух одинаковых автоматах к концу дня может закончиться кофе. Вероятность того, что кофе закончится в первом автомате, равна $0{,}3$, во втором --- также $0{,}3$. Вероятность, что кофе останется в обоих автоматах, равна $0{,}55$. Найдите вероятность того, что кофе закончится в обоих.
\field{Ваше решение и ответ}
На листе отсутствуют.
\field{Где возникает затруднение}
Письменных шагов нет. Важно различать события «хотя бы в одном» и «в обоих». Независимость автоматов в условии не утверждается.
\field{Ключевая идея}
\keyidea Для событий $A$ и $B$ справедливо
\[P(A\cup B)=P(A)+P(B)-P(A\cap B).\]
\field{Полное правильное решение}
Обозначим через $A$ событие «кофе закончился в первом автомате», через $B$ --- «во втором».
\[P(A)=0{,}3,\qquad P(B)=0{,}3.\]
Вероятность того, что кофе закончится хотя бы в одном автомате, равна вероятности, дополнительной к событию «кофе остался в обоих»:
\[P(A\cup B)=1-0{,}55=0{,}45.\]
Теперь выразим вероятность пересечения:
\begin{align*}
P(A\cap B)&=P(A)+P(B)-P(A\cup B)\\
&=0{,}3+0{,}3-0{,}45=0{,}15.
\end{align*}
\field{Проверка}
$0{,}15$ не превышает ни $0{,}3$, ни $0{,}3$. Событие «кофе закончился в обоих автоматах» действительно не может быть вероятнее каждого из событий по отдельности.
\result Вероятность равна $0{,}15$.
\field{Задача на закрепление}
Вероятность отказа первого устройства равна $0{,}40$, второго --- $0{,}35$. Вероятность, что оба устройства работают, равна $0{,}45$. Найдите вероятность отказа обоих устройств.

\taskstart{6}{Теория вероятностей · закон распределения}
\field{Условие}
Случайная величина $Y$ задана таблицей:
\[\begin{array}{c|rrrrrr}Y&0&2&4&8&16&32\\ \hline P&0{,}10&0{,}20&0{,}25&0{,}25&0{,}15&0{,}05\end{array}\]
Найдите $P(-2<Y<8)$.
\field{Ваше решение и ответ}
На листе отсутствуют.
\field{Где возникает затруднение}
Решение не записано. Требуется отобрать ровно те значения величины, которые удовлетворяют двум строгим неравенствам одновременно.
\field{Ключевая идея}
\keyidea Для дискретной случайной величины вероятность события равна сумме вероятностей подходящих значений.
\field{Полное правильное решение}
В интервал $(-2;8)$ из таблицы входят только $0$, $2$ и $4$:
\[-2<0<8,\qquad -2<2<8,\qquad -2<4<8.\]
Число $8$ исключается, поскольку верхняя граница строгая. Тогда
\[P(-2<Y<8)=P(Y=0)+P(Y=2)+P(Y=4)\]
\[=0{,}10+0{,}20+0{,}25=0{,}55.\]
\field{Проверка}
Сумма всех вероятностей равна $1$. Для значений вне события получаем $0{,}25+0{,}15+0{,}05=0{,}45$. Значит, искомая вероятность $1-0{,}45=0{,}55$.
\result $P(-2<Y<8)=0{,}55$.
\field{Задача на закрепление}
Случайная величина $X$ принимает значения $-3$, $0$, $2$, $5$, $9$ с вероятностями $0{,}12$, $0{,}18$, $0{,}25$, $0{,}30$, $0{,}15$ соответственно. Найдите $P(-1<X<9)$.

\taskstart{9}{Исследование функции по графику производной}
\field{Условие}
По данному графику производной $y=f'(x)$ определите количество точек максимума функции $f(x)$ на отрезке $[-18;1]$.
\field{Ваше решение и ответ}
Записан ответ $5$; промежуточные рассуждения отсутствуют.
\field{Ошибка и её локализация}
\wrong Итоговое число точек максимума неверно. По одному записанному ответу нельзя надёжно выяснить, какие именно точки были посчитаны. Проверяем сам критерий максимума.
\field{Ключевая идея}
\keyidea Точка максимума $f(x)$ возникает там, где $f'(x)$ меняет знак с «плюса» на «минус» при увеличении $x$. Вершины графика $f'(x)$ сами по себе такого вывода не дают.
\field{Полное правильное решение}
\begin{enumerate}
\item Рассматриваем только часть графика производной при $-18\le x\le1$.
\item Отмечаем точки пересечения графика $f'(x)$ с осью абсцисс.
\item Для каждой точки сравниваем знак производной непосредственно слева и справа. Переход $-\to+$ даёт минимум функции, переход $+\to-$ --- максимум.
\item На заданном отрезке видны ровно \textbf{два} перехода $+\to-$. Поэтому точек максимума две.
\end{enumerate}
\field{Проверка}
Переходы графика производной через ось снизу вверх не учитываются как максимумы. Также не учитывается переход $+\to-$, расположенный левее $-18$.
\result Количество точек максимума: $\boxed{2}$.
\field{Задача на закрепление}
На четырёх последовательных промежутках знаки производной $f'(x)$ равны соответственно $+$, $-$, $+$, $-$. В точках раздела производная обращается в нуль. Сколько точек максимума функции находятся внутри рассматриваемого интервала?

\taskstart{11}{Задача на движение · катер и течение}
\begingroup\small\setlength{\parskip}{4pt}
\field{Условие}
Катер движется по течению реки, затем после стоянки возвращается в исходный пункт. Собственная скорость катера $28$ км/ч, скорость течения $2$ км/ч; стоянка длится $4$ часа, возвращение происходит через $18$ часов после отплытия. Найдите полный пройденный путь.
\field{Ваше решение}
В таблице указаны скорости $30$ и $26$ км/ч, расстояния по $x$, но время записано как $30/x$ и $26/x$. Далее составлено уравнение
\[\frac{30}{x}+\frac{26}{x}+4=18,\]
получено $x=4$ и ответ $8$ км.
\field{Где именно возникает ошибка}
\wrong Скорости $30=28+2$ и $26=28-2$ найдены верно. Первый неверный шаг --- заполнение графы времени: вместо $x/30$ и $x/26$ написаны обратные отношения.
\field{Почему это неверно}
Из формулы пути $S=vt$ следует $t=S/v$. Деление скорости на расстояние даёт размерность $1/\text{ч}$, а время должно измеряться в часах. Поэтому запись $v/S$ здесь непригодна.
\field{Ключевая идея}
\keyidea Составлять таблицу «путь --- скорость --- время» и в столбце времени неизменно использовать $t=S/v$.
\field{Исправление последнего верного шага}
Сохраняем найденные скорости: по течению --- $30$ км/ч, против течения --- $26$ км/ч. Пусть расстояние в одну сторону равно $x$ км. Тогда
\[t_1=\frac{x}{30},\qquad t_2=\frac{x}{26}.\]
\field{Полное правильное решение}
Суммарное время движения и стоянки составляет $18$ часов:
\begin{align*}
\frac{x}{30}+\frac{x}{26}+4&=18,\\
\frac{x}{30}+\frac{x}{26}&=14,\\
\frac{26x+30x}{780}&=14,\\
\frac{56x}{780}&=14,\\
56x&=10920,\\
x&=195.
\end{align*}
Это расстояние от исходного пункта до места стоянки. За весь рейс катер проходит $2x=390$ км.
\field{Проверка}
По течению катер двигался $195/30=6{,}5$ ч, против течения --- $195/26=7{,}5$ ч. С учётом стоянки $6{,}5+7{,}5+4=18$ ч.
\result Полный путь --- $\boxed{390\text{ км}}$.
\field{Задача на закрепление}
Собственная скорость катера $20$ км/ч, скорость течения $4$ км/ч. Катер прошёл по течению до пристани, простоял там $2$ часа и вернулся в исходную точку через $12$ часов после отплытия. Найдите полный путь катера.

\endgroup
\taskstart{13}{Финансовая задача · кредит с ежегодным изменением долга}
\field{Условие}
В июле 2031 года Леонид берёт кредит $1{,}2$ млн рублей на три года. В январе каждого года долг увеличивается на $15\%$ от долга на конец предыдущего года. С февраля по июнь выплачивается часть долга. В июле 2032 и 2033 годов долг должен быть на $30\%$ меньше, чем в июле предыдущего года. К июлю 2034 года кредит погашен полностью. Найдите общую сумму выплат.
\field{Ваше решение и ответ}
На листе отсутствуют.
\field{Где возникает затруднение}
Письменных вычислений нет. Для решения важно различать остаток долга в июле и долг после январского увеличения.
\field{Ключевая идея}
\keyidea Каждый год сначала умножаем июльский остаток предыдущего года на $1{,}15$, затем вычитаем плановый июльский остаток текущего года.
\field{Полное правильное решение}
В июле 2031 года долг равен $1\,200\,000$ руб.
\[D_{2032}=1\,200\,000\cdot0{,}7=840\,000\text{ руб.}\]
\[D_{2033}=840\,000\cdot0{,}7=588\,000\text{ руб.}\]
\[D_{2034}=0.\]
В 2032 году после начисления процентов долг составляет
\[1\,200\,000\cdot1{,}15=1\,380\,000.\]
Чтобы к июлю осталось $840\,000$, необходимо выплатить
\[P_{2032}=1\,380\,000-840\,000=540\,000.\]
В 2033 году аналогично:
\[P_{2033}=840\,000\cdot1{,}15-588\,000=378\,000.\]
В 2034 году погашаем весь долг после начисления процентов:
\[P_{2034}=588\,000\cdot1{,}15=676\,200.\]
Общая сумма платежей:
\[540\,000+378\,000+676\,200=1\,594\,200\text{ руб.}\]
\field{Проверка}
Переплата составляет $1\,594\,200-1\,200\,000=394\,200$ руб. Она равна сумме начисленных процентов: $180\,000+126\,000+88\,200=394\,200$ руб.
\result Общая сумма выплат --- $\boxed{1\,594\,200\text{ руб.}}$.
\field{Задача на закрепление}
В июле взяли кредит $900\,000$ рублей на три года. Каждый январь остаток долга увеличивается на $20\%$, а к июлю первого и второго годов выплат остаток долга должен уменьшаться на $20\%$ относительно предыдущего июля. К июлю третьего года кредит погашен. Найдите общую сумму выплат.

\taskstart{14}{Тригонометрическое уравнение · формулы приведения и ОДЗ}
\field{Условие}
Решите уравнение
\[\sqrt{2\cos^3x-\sin^2x-2\cos x-\sin x}
=\sqrt{\cos\left(\frac\pi2+x\right)}.\]
Найдите все корни на отрезке $\left[-4\pi;-\frac{5\pi}{2}\right]$.
\field{Ваше решение}
В первой строке преобразований записано
\[\cos\left(\frac\pi2+x\right)=\sin x.\]
Затем из равенства квадратных корней выполнен переход к равенству, в котором правая часть равна $\sin x$. В конце получена ветвь $\sin x=0$, но решение второй ветви и отбор корней не завершены.
\field{Где именно возникает ошибка}
\wrong Первый неверный шаг --- формула приведения: потерян знак минус. Правильно
\[\cos\left(\frac\pi2+x\right)=-\sin x.\]
Кроме того, из $\sqrt{A}=\sqrt{B}$ следует $A=B$ при $A\ge0$, $B\ge0$; вообще говоря, $\sqrt{B}\ne B$.
\field{Почему это важно}
Неверный знак меняет уравнение и допустимые значения $x$. Подкоренное выражение справа должно быть неотрицательным, поэтому $-\sin x\ge0$, то есть $\sin x\le0$.
\field{Ключевая идея}
\keyidea Сначала применить формулу приведения, записать ОДЗ, затем приравнять подкоренные выражения и разложить полученную разность на множители.
\field{Полное правильное решение}
Шаг 1. Приводим правую часть:
\[\cos\left(\frac\pi2+x\right)=-\sin x.\]
ОДЗ требует $\sin x\le0$ и неотрицательности левого подкоренного выражения.

Шаг 2. Приравниваем подкоренные выражения:
\[2\cos^3x-\sin^2x-2\cos x-\sin x=-\sin x.\]
Слагаемые $-\sin x$ сокращаются:
\[2\cos^3x-2\cos x-\sin^2x=0.\]
Шаг 3. Выносим общий множитель и применяем основное тригонометрическое тождество:
\begin{align*}
2\cos x(\cos^2x-1)-\sin^2x&=0,\\
-2\cos x\sin^2x-\sin^2x&=0,\\
-\sin^2x(2\cos x+1)&=0.
\end{align*}
Отсюда две ветви:
\[\sin x=0\quad\text{или}\quad\cos x=-\frac12.\]
Шаг 4. Решаем и проверяем ОДЗ.

Из $\sin x=0$ получаем $x=\pi n$, $n\in\mathbb Z$. Для этих значений обе части исходного равенства равны нулю.

Из $\cos x=-\frac12$ получаем $x=\frac{2\pi}{3}+2\pi n$ или $x=\frac{4\pi}{3}+2\pi n$. Первое семейство имеет $\sin x>0$ и не удовлетворяет ОДЗ; второе допустимо.

Значит, все решения:
\[x=\pi n\quad\text{или}\quad x=\frac{4\pi}{3}+2\pi n,\qquad n\in\mathbb Z.\]
Шаг 5. Отбираем корни на данном отрезке:
\begin{align*}
\pi n &\in\left[-4\pi;-\frac{5\pi}{2}\right]
&&\Longrightarrow n=-4,-3;\\
\frac{4\pi}{3}+2\pi n &\in\left[-4\pi;-\frac{5\pi}{2}\right]
&&\Longrightarrow n=-2.
\end{align*}
\field{Проверка}
Для $x=-4\pi$ и $x=-3\pi$ обе части равны $0$. Для $x=-8\pi/3$ выполнено $\cos x=-1/2$, $\sin x=-\sqrt3/2$; подкоренные выражения совпадают и равны $\sqrt3/2>0$.
\result На заданном отрезке: $\boxed{-4\pi,\ -3\pi,\ -\frac{8\pi}{3}}$.
\field{Задача на закрепление}
Решите уравнение
\[\sqrt{2\cos^3x-\sin^2x-2\cos x+\sin x}
=\sqrt{\cos\left(\frac\pi2-x\right)}\]
и найдите все его корни на отрезке $[0;2\pi]$.
\end{document}